% Compile with XeLaTeX (recommended) or pdfLaTeX. Requires the 'kotex' package
% (bundled in TeX Live / available on Overleaf). On Overleaf: Menu > Compiler > XeLaTeX.
\documentclass[11pt]{article}
\usepackage[margin=1in]{geometry}
\usepackage{kotex}
\usepackage{parskip}
\usepackage{titlesec}
\usepackage{xcolor}
\definecolor{accent}{HTML}{1D4ED8}
\titleformat{\section}{\normalfont\large\bfseries\color{accent}}{}{0em}{}
\titlespacing{\section}{0pt}{1.2em}{0.4em}

\title{\textbf{From Proximity to Tactical Advantage}\\[2pt]
\large 발표 대본 --- 슬라이드별 (한국어)}
\author{백성은 \,\textbar\, 부산대학교 \,\textbar\, 그래프데이터분석}
\date{}

\begin{document}
\maketitle
\noindent\textit{15장 덱(\texttt{graph\_deck.pdf})용 발표 대본입니다. 구어체(존댓말), 약 9{--}11분 분량. 핵심 발견을 앞에서부터 제시하는 서사입니다 --- 평범한 상호작용 위상은 평활화(smoothing)처럼 작동했고, 이것이 관계 자체를 재설계해야 한다는 동기를 줍니다.}
\vspace{0.5em}\hrule

\section{슬라이드 1 --- 표지}
안녕하세요. 부산대학교 백성은입니다. 오늘은 그래프 모델로 배틀그라운드의 생존을 예측한 연구를 발표하겠습니다. 출발점은 단순한 질문이었습니다 --- 팀들을 ``상호작용''으로 연결한 그래프 구조가 정말로 누가 살아남을지 예측하는 데 도움이 될까요? 미리 말씀드리면 결과는 의외였습니다. 흔히 떠올리는 상호작용 엣지는 대부분 ``평활화''처럼 작동했을 뿐, 실제 전술적 신호는 아니었습니다. 그리고 바로 이 부정적 결과가 이 연구를 흥미롭게 만듭니다.

\section{슬라이드 2 --- 수축하는 서식지 속 팀 관계로서의 PUBG}
먼저 관점입니다. 배틀그라운드 한 판은 고립된 개인들이 아니라, 점점 줄어드는 안전구역 안에서 움직이고 싸우는 ``팀들''입니다. 이건 Erangel 맵의 실제 한 경기입니다. 점 하나가 한 팀이고, 선 하나가 우리 모델이 인코딩하는 관계입니다 --- 누가 가까이 있는가, 누가 누구를 쏘는가, 누가 존 압박에 함께 몰리는가. 우승팀은 그림 설명용으로만 표시했고 모델 입력으로는 절대 쓰이지 않습니다. 요점은, 여기서 생존은 본질적으로 ``관계적''이고 ``공간적''이라는 것입니다.

\section{슬라이드 3 --- 목적: 관계적, 팀 단위 관점}
그러면 선행연구와 무엇이 다를까요? 기존 PUBG 연구는 대부분 플레이어 한 명씩을 봅니다 --- 공격성을 점수화하고, 플레이 스타일로 분류하고, 그 플레이어의 생존을 예측합니다. 우리는 대신 한 경기를 ``수축하는 서식지 안에서 진화하는 팀 관계의 그래프''로 봅니다. 목적은 phase별로 어느 팀이 다음에 탈락하는지, 그리고 각 팀이 최종 몇 등인지 예측하는 것 --- 그리고 무엇보다, 그 관계 구조 자체가 신호를 담고 있는지 검증하는 것입니다.

\section{슬라이드 4 --- 선행연구}
가장 가까운 선행연구는 2025년 Sagong 등의 연구입니다. 이들은 phase마다 FightScore와 ActiveScore를 만들어 플레이어를 네 유형으로 군집화하고 트리 모델로 예측합니다 --- 하지만 ``솔로'' 플레이어만 대상으로 하고, 팀은 의도적으로 제외했습니다. 우리는 바로 그들이 제외한 ``팀 관계'' 차원을 다시 가져오고, 단순히 성능이 잘 나오는지가 아니라 ``null-control 하에서 그 위상이 실제로 신호를 담는지''를 신중히 검증합니다.

\section{슬라이드 5 --- 데이터와 예측 과제}
데이터는 매치 텔레메트리에서 만든 squad 시공간 데이터베이스입니다 --- 9개 맵에서 약 1만 5천 경기를 경기 단위로 학습/검증/테스트로 나눴고, 라벨이 붙은 팀-phase 셀이 약 86만 개입니다. 이제 라벨링인데, 여기선 누수가 생기기 쉽습니다. 모든 예측은 ``인과적 컷''에서 이뤄집니다 --- 피처는 phase $p$까지만 관측하고 $p\!+\!1$을 예측합니다. 팀이 살아있는 동안에만 채점합니다 --- 이게 at-risk 마스크고, 이미 탈락한 팀은 예측하지 않습니다. 두 타깃은 ``다음 phase에 탈락하는가''(AUC로 채점)와 ``최종 몇 등인가''(Spearman으로 채점)입니다. 그리고 이름 기반 차단 목록으로 어떤 결과(outcome) 필드도 입력에 들어가지 못하게 막습니다.

\section{슬라이드 6 --- 동적 이종(heterogeneous) 그래프}
구체적으로, 한 경기는 ``동적 이종 그래프''입니다. ``이종''인 이유는 팀들이 세 가지 서로 다른 관계 유형으로 연결되고 이를 별도 채널로 유지하기 때문입니다 --- proximity, combat, co-pressure. 그리고 R-GCN이 각 관계를 자신만의 행렬로 처리합니다. ``동적''인 이유는 그래프가 phase마다 다시 구성되고 10초 스냅샷으로도 샘플링되며, causal temporal encoder가 그 진화를 읽기 때문입니다.

\section{슬라이드 7 --- 세 관계는 어떻게 만드는가}
세 관계는 서로 다르게 정의됩니다. proximity는 대칭적인 거리 커널입니다 --- 가까운 팀일수록 강한 엣지. combat은 방향성이 있습니다 --- 공격자에서 피해자로, 피해량/기절/킬로 가중됩니다. co-pressure는 대칭이지만 ``게이트''가 걸립니다 --- 가까운 두 팀이 ``둘 다'' 존 압박을 받을 때만 활성화됩니다. 각 채널은 관계별로 행 정규화됩니다.

\section{슬라이드 8 --- 아키텍처}
아키텍처입니다. 입력은 노드 피처와 세 관계 채널을 가진 팀 그래프입니다. 핵심은 시공간 블록으로 두 번 적용됩니다 --- 각 팀이 관계별로 이웃의 메시지를 집계하는 그래프 레이어와, phase에 걸친 causal temporal convolution입니다. 그 결과 팀마다 64차원 임베딩이 나오고, 이게 두 헤드로 들어갑니다 --- 탈락을 위한 hazard MLP, 그리고 최종 순위를 위한 디코더입니다. 모델 변종 사이에서는 그래프 인코더만 교체됩니다.

\section{슬라이드 9 --- 탈락 예측 결과}
결과입니다. 탈락 예측에서는 모든 그래프 인코더 --- GCN, GraphSAGE, GAT, R-GCN --- 가 서로 AUC 0.02 안쪽에서 거의 붙어 있고, no-graph 베이스라인보다 살짝 위입니다. 인코더들이 이렇게 촘촘히 모여 있다는 건, 이미 ``아키텍처가 핵심 지렛대가 아니다''라는 것을 말해 줍니다.

\section{슬라이드 10 --- 순위 예측 결과}
최종 순위에서는 상관이 phase 3에서 6으로 분명히 올라갑니다. 그런데 --- 이게 중요합니다 --- ``그래프가 전혀 없을 때도'' 똑같은 상승이 나타납니다. 경기가 진행될수록 at-risk 집합이 약 7배 줄어들어, 남은 소수의 팀은 그냥 순서를 매기기가 쉬워집니다. 그래프는 곡선 전체를 조금 위로 올릴 뿐 그 상승 자체를 만들지는 않습니다. 모델의 이득이 아니라 ``survivor 선택 효과''입니다.

\section{슬라이드 11 --- 성능 종합}
성능을 요약하면 이렇습니다 --- 그래프를 넣으면 수치가 조금 오르고, 인코더 간 차이는 아주 작습니다. 하지만 ``성능 향상''이 곧 ``위상 신호''는 아닙니다 --- 그리고 방금 봤듯 순위 이득은 survivor 효과입니다. 이것이 다음 슬라이드의 핵심 검증으로 이어집니다.

\section{슬라이드 12 --- Null-control: 핵심 결과}
이게 핵심 결과입니다. 관측된 위상을 ``edge-shuffle'' null과 비교합니다. edge-shuffle은 각 관계의 가중치 분포와 밀도는 그대로 두되, 끝점을 살아있는 팀들 사이에서 무작위로 재배치합니다 --- 즉 다른 건 다 보존하고 ``실제 연결(wiring)''만 파괴합니다. 실제 위상이 의미가 있다면 관측이 shuffle을 이겨야 합니다. 그렇지 않았습니다 --- 모든 시드에서 shuffle이 같거나 더 좋았고, 초반 phase에서는 유의하게 더 좋았습니다. 정직한 결론은, 관측된 위상이 단순 그래프 평균화 이상의 ``쓸 만한 신호''를 더하지 않는다는 것입니다.

\section{슬라이드 13 --- 관계의 재설계 (예비)}
왜 실패했을까요? proximity, combat, co-pressure는 상호작용이 ``존재하는지''와 ``얼마나 강한지''를 인코딩할 뿐, ``누가 유리한지''를 담지 않기 때문입니다. 예비적 후속으로, 우리는 엣지가 ``상대적 우위''를 담도록 재설계했습니다 --- 고도, 인원수, 체력, 존 위치의 차이를, 접촉(contact) 시에만 게이트해서요. 이 관계는 null-control 대비 작지만 일관된 양의 신호를 보였습니다. 효과는 작고 아직 초기 결과지만 방향은 맞습니다.

\section{슬라이드 14 --- 한계점과 기여점}
한계를 정직하게 말씀드리면 --- 관측 위상은 본질적으로 평활화이고, 맵 그림은 한 경기 예시이며, advantage 엣지는 예비 결과입니다. 기여는 PUBG 생존을 위한 phase-causal/팀 단위 그래프 정식화, 그 평활화를 드러낸 null-control 프로토콜, 관계 설계 가설, 그리고 스냅샷/phase/hazard/rank를 잇는 다해상도 파이프라인입니다.

\section{슬라이드 15 --- 결론}
결론은 이렇습니다 --- 위상이 도움이 되는지를 결정하는 것은 GNN 아키텍처가 아니라 ``관계의 유형''입니다. 팀을 단순히 거리나 교전으로 잇는 것은 상태를 평균낼 뿐입니다. 가치는 게임 메커닉에 맞는 관계 --- 상대적 전술 우위 --- 를 설계하고, GNN은 그것을 ``읽는'' 도구로만 쓰는 데 있습니다. 감사합니다. 질문 받겠습니다.

\vspace{1em}\hrule
\vspace{0.4em}
\noindent\small\textit{예상 질문 대비: (1) 여기서 ``causal''은 인과추론이 아니라 누수 없는(leakage-safe)/phase-causal 컷을 뜻합니다. (2) N=128은 패딩 크기이고 실제 squad 팀 수는 평균 약 27팀입니다. T=8은 phase 슬롯, 채점되는 cut-phase는 7개(P1{--}P7)입니다. (3) no-graph 베이스라인은 그래프 파생 노드 피처(network-role)를 제외한 core\_habitat 피처셋을 쓰므로 null 비교가 오염되지 않았습니다. (4) edge-shuffle은 match-phase별 가중치 분포/밀도는 보존하고 끝점만 무작위화합니다. (5) 순위 상승은 survivor 집합 축소 효과로, 그래프 없이도 동일하게 나타납니다. (6) advantage 엣지의 $+0.003$ AUC는 작은, 가설-생성 수준의 결과이며 주요 주장이 아닙니다.}

\end{document}
